%! Direct, Inverse \& Joint Variation — Unit 4, Lesson 4.7 — Homework
\documentclass[10pt]{article}
\usepackage{algebra2-article}
\usepackage{algebra2-boxes}
\usepackage{grading}
\gradingassignment{alg2-u04-l07-hw}
\gradingcourse{algebra_2}
\usepackage{algebra2-graded}
\newcommand{\TallMath}[1]{$\displaystyle #1\rule[-1.4em]{0pt}{3.2em}$}
% Standard coordinate grid + axes: {xmin}{xmax}{ymin}{ymax}
\newcommand{\lgrid}[4]{%
  \draw[step=1, linegray!40, very thin] (#1,#3) grid (#2,#4);
  \draw[->, semithick, charcoal] (#1,0)--(#2,0) node[below right, font=\scriptsize]{$x$};
  \draw[->, semithick, charcoal] (0,#3)--(0,#4) node[left, font=\scriptsize]{$y$};}

\begin{document}

\pageheader{Unit 4, Lesson 4.7}{Homework}

\vspace{0.10in}

\begin{notesbox}{Practice}
Every table gets \emph{both} tests: the quotients $y\div x$ and the products $x\cdot y$. Whichever row
is constant names the variation --- and if neither is, the answer is ``neither.''

\begin{enumerate}[leftmargin=1.9em, itemsep=10pt]
  \item \textbf{Name the variation} and give $k$ when there is one.
    \begin{center}
    \begin{minipage}{0.47\linewidth}\centering
    \renewcommand{\arraystretch}{1.2}
    \begin{tabular}{|l|c|c|c|c|}
    \hline
    \rowcolor{forestbg}
    $x$ & $4$ & $6$ & $9$ & $12$ \\ \hline
    $y$ & $10$ & $15$ & $22.5$ & $30$ \\ \hline
    \end{tabular}\\[2pt]
    {\small (a) \inblank{\blank{2.2cm}}{\gradedblank[id=q1a, type=short, points=1, answer={direct}, check=text]{2.2cm-2pt}{10pt}} \; $k=\inblank{\blank{1.2cm}}{\gradedblank[id=q1b, type=short, points=1, answer={2.5}, check=numeric]{1.2cm-2pt}{10pt}}$}
    \end{minipage}\hfill
    \begin{minipage}{0.47\linewidth}\centering
    \renewcommand{\arraystretch}{1.2}
    \begin{tabular}{|l|c|c|c|c|}
    \hline
    \rowcolor{forestbg}
    $x$ & $3$ & $4$ & $6$ & $12$ \\ \hline
    $y$ & $8$ & $6$ & $4$ & $2$ \\ \hline
    \end{tabular}\\[2pt]
    {\small (b) \inblank{\blank{2.2cm}}{\gradedblank[id=q1c, type=short, points=1, answer={inverse}, check=text]{2.2cm-2pt}{10pt}} \; $k=\inblank{\blank{1.2cm}}{\gradedblank[id=q1d, type=short, points=1, answer={24}, check=numeric]{1.2cm-2pt}{10pt}}$}
    \end{minipage}
    \end{center}
    \vspace{3pt}
    \begin{center}
    \begin{minipage}{0.47\linewidth}\centering
    \renewcommand{\arraystretch}{1.2}
    \begin{tabular}{|l|c|c|c|c|}
    \hline
    \rowcolor{forestbg}
    $x$ & $1$ & $3$ & $5$ & $7$ \\ \hline
    $y$ & $4$ & $10$ & $16$ & $22$ \\ \hline
    \end{tabular}\\[2pt]
    {\small (c) \inblank{\blank{2.2cm}}{\gradedblank[id=q1e, type=short, points=1, answer={neither}, check=text]{2.2cm-2pt}{10pt}} \; $k=\inblank{\blank{1.2cm}}{\gradedblank[id=q1f, type=short, points=1, answer={none}, check=text]{1.2cm-2pt}{10pt}}$}
    \end{minipage}\hfill
    \begin{minipage}{0.47\linewidth}\centering
    \renewcommand{\arraystretch}{1.2}
    \begin{tabular}{|l|c|c|c|c|}
    \hline
    \rowcolor{forestbg}
    $x$ & $3$ & $6$ & $7$ & $10$ \\ \hline
    $y$ & $21$ & $42$ & $49$ & $70$ \\ \hline
    \end{tabular}\\[2pt]
    {\small (d) \inblank{\blank{2.2cm}}{\gradedblank[id=q1g, type=short, points=1, answer={direct}, check=text]{2.2cm-2pt}{10pt}} \; $k=\inblank{\blank{1.2cm}}{\gradedblank[id=q1h, type=short, points=1, answer={7}, check=numeric]{1.2cm-2pt}{10pt}}$}
    \end{minipage}
    \end{center}

  \item \textbf{Find $k$, write the equation, answer the question.}
    \begin{enumerate}[label=(\alph*), leftmargin=2.1em, itemsep=5pt]
      \item $y$ varies \emph{directly} as $x$; $y=24$ when $x=8$. \hfill $k=\inblank{\blank{1.2cm}}{\gradedblank[id=q2a, type=short, points=1, answer={3}, check=numeric]{1.2cm-2pt}{10pt}}$ \;
            $y=\inblank{\blank{1.6cm}}{\gradedblank[id=q2b, type=short, points=1, answer={y=3x}, display={$3x$}]{1.6cm-2pt}{10pt}}$ \; Find $y$ when $x=15$: \inblank{\blank{1.4cm}}{\gradedblank[id=q2c, type=short, points=1, answer={45}, check=numeric]{1.4cm-2pt}{10pt}}
      \item $y$ varies \emph{inversely} as $x$; $y=7$ when $x=6$. \hfill $k=\inblank{\blank{1.2cm}}{\gradedblank[id=q2d, type=short, points=1, answer={42}, check=numeric]{1.2cm-2pt}{10pt}}$ \;
            $y=\inblank{\blank{1.6cm}}{\gradedblank[id=q2e, type=short, points=1, answer={y=42/x}, display={$42/x$}]{1.6cm-2pt}{10pt}}$ \; Find $y$ when $x=14$: \inblank{\blank{1.4cm}}{\gradedblank[id=q2f, type=short, points=1, answer={3}, check=numeric]{1.4cm-2pt}{10pt}}
      \item $y$ varies \emph{directly} as $x$; $y=4.5$ when $x=6$. \hfill $k=\inblank{\blank{1.2cm}}{\gradedblank[id=q2g, type=short, points=1, answer={0.75}, check=numeric]{1.2cm-2pt}{10pt}}$ \;
            $y=\inblank{\blank{1.6cm}}{\gradedblank[id=q2h, type=short, points=1, answer={y=0.75x}, display={$0.75x$}]{1.6cm-2pt}{10pt}}$ \; Find $x$ when $y=12$: \inblank{\blank{1.4cm}}{\gradedblank[id=q2i, type=short, points=1, answer={16}, check=numeric]{1.4cm-2pt}{10pt}}
      \item $y$ varies \emph{inversely} as $x$; $y=2.5$ when $x=16$. \hfill $k=\inblank{\blank{1.2cm}}{\gradedblank[id=q2j, type=short, points=1, answer={40}, check=numeric]{1.2cm-2pt}{10pt}}$ \;
            $y=\inblank{\blank{1.6cm}}{\gradedblank[id=q2k, type=short, points=1, answer={y=40/x}, display={$40/x$}]{1.6cm-2pt}{10pt}}$ \; Find $x$ when $y=8$: \inblank{\blank{1.4cm}}{\gradedblank[id=q2l, type=short, points=1, answer={5}, check=numeric]{1.4cm-2pt}{10pt}}
      \item $y$ varies \emph{inversely} as $x$; $y=20$ when $x=3$. \hfill $k=\inblank{\blank{1.2cm}}{\gradedblank[id=q2m, type=short, points=1, answer={60}, check=numeric]{1.2cm-2pt}{10pt}}$ \;
            $y=\inblank{\blank{1.6cm}}{\gradedblank[id=q2n, type=short, points=1, answer={y=60/x}, display={$60/x$}]{1.6cm-2pt}{10pt}}$ \; Find $y$ when $x=5$: \inblank{\blank{1.4cm}}{\gradedblank[id=q2o, type=short, points=1, answer={12}, check=numeric]{1.4cm-2pt}{10pt}}
      \item $y$ varies \emph{directly} as $x$; $y=-14$ when $x=4$. \hfill $k=\inblank{\blank{1.2cm}}{\gradedblank[id=q2p, type=short, points=1, answer={-3.5}, check=numeric]{1.2cm-2pt}{10pt}}$ \;
            $y=\inblank{\blank{1.6cm}}{\gradedblank[id=q2q, type=short, points=1, answer={y=-3.5x}, display={$-3.5x$}]{1.6cm-2pt}{10pt}}$ \; Find $y$ when $x=10$: \inblank{\blank{1.4cm}}{\gradedblank[id=q2r, type=short, points=1, answer={-35}, check=numeric]{1.4cm-2pt}{10pt}}
    \end{enumerate}
    \vspace{2pt}
    {\small Part (f) has a negative $k$. Nothing breaks --- but say what it does to the graph:
    \inblank{\blank{5.0cm}}{\gradedblank[id=q2s, type=rubric, points=2, display={the line still passes through the origin, but it slopes downward, so $y$ decreases as $x$
    increases}]{5.0cm-2pt}{10pt}}}

  \item \textbf{No table given. Decide from the situation, and say why.}
    \begin{enumerate}[label=(\alph*), leftmargin=2.1em, itemsep=5pt]
      \item The perimeter $P$ of a square and its side length $s$. \; \inblank{\blank{2.0cm}}{\gradedblank[id=q3a, type=short, points=1, answer={direct}, check=text]{2.0cm-2pt}{10pt}}, because
            $P=\inblank{\blank{1.4cm}}{\gradedblank[id=q3b, type=short, points=1, answer={P=4s}, display={$4s$}]{1.4cm-2pt}{10pt}}$, so $k=\inblank{\blank{1.0cm}}{\gradedblank[id=q3c, type=short, points=1, answer={4}, check=numeric]{1.0cm-2pt}{10pt}}$.
      \item One pizza cut into $8$ slices is shared equally among $n$ people. The number of slices
            each person gets and $n$: \; \inblank{\blank{2.0cm}}{\gradedblank[id=q3d, type=short, points=1, answer={inverse}, check=text]{2.0cm-2pt}{10pt}}, $k=\inblank{\blank{1.0cm}}{\gradedblank[id=q3e, type=short, points=1, answer={8}, check=numeric]{1.0cm-2pt}{10pt}}$.
      \item A taxi charges $\$3.00$ just to start the meter, plus $\$2.00$ per mile. Is the fare
            directly proportional to the miles driven? \inblank{\blank{1.2cm}}{\gradedblank[id=q3f, type=short, points=1, answer={no}, check=text]{1.2cm-2pt}{10pt}} \\
            Why not? \inblank{\blank{6.0cm}}{\gradedblank[id=q3g, type=rubric, points=2, display={the fare is $2m+3$: zero miles still costs $\$3$, so the graph misses the
            origin and fare$\,\div\,$miles is not constant}]{6.0cm-2pt}{10pt}}
    \end{enumerate}

  \item \textbf{Read the graphs.} One of these is a direct variation and one is an inverse variation.
    \begin{center}
    \begin{minipage}{0.46\linewidth}\centering
    \begin{tikzpicture}[scale=0.30]
      \lgrid{-7}{7}{-9}{9}
      \begin{scope}
        \clip (-7,-9) rectangle (7,9);
        \draw[very thick, forest] (-6,-9)--(6,9);
      \end{scope}
      \fill[charcoal] (4,6) circle (3.6pt) node[above left, font=\scriptsize]{$(4,6)$};
    \end{tikzpicture}\\[2pt]
    {\small \textbf{Graph I}}
    \end{minipage}\hfill
    \begin{minipage}{0.46\linewidth}\centering
    \begin{tikzpicture}[scale=0.30]
      \lgrid{-9}{9}{-9}{9}
      \begin{scope}
        \clip (-9,-9) rectangle (9,9);
        \draw[very thick, forest, domain=1.34:9, samples=150, smooth] plot (\x,{12/(\x)});
        \draw[very thick, forest, domain=-9:-1.34, samples=150, smooth] plot (\x,{12/(\x)});
      \end{scope}
      \fill[charcoal] (2,6) circle (3.6pt) node[above right, font=\scriptsize]{$(2,6)$};
      \fill[charcoal] (6,2) circle (3.6pt);
    \end{tikzpicture}\\[2pt]
    {\small \textbf{Graph II}}
    \end{minipage}
    \end{center}
    \vspace{-2pt}
    \begin{enumerate}[label=(\alph*), leftmargin=2.1em, itemsep=5pt]
      \item Graph I is a \inblank{\blank{2.0cm}}{\gradedblank[id=q4a, type=short, points=1, answer={direct}, check=text]{2.0cm-2pt}{10pt}} variation; the giveaway is that it is a line through the
            \inblank{\blank{1.8cm}}{\gradedblank[id=q4b, type=short, points=1, answer={origin}, check=text]{1.8cm-2pt}{10pt}}. Using the marked point, $k=\inblank{\blank{1.0cm}}{\gradedblank[id=q4c, type=short, points=1, answer={1.5}, check=numeric]{1.0cm-2pt}{10pt}}$ and $y=\inblank{\blank{1.4cm}}{\gradedblank[id=q4d, type=short, points=1, answer={y=1.5x}, display={$1.5x$}]{1.4cm-2pt}{10pt}}$.
      \item Graph II is a \inblank{\blank{2.0cm}}{\gradedblank[id=q4e, type=short, points=1, answer={inverse}, check=text]{2.0cm-2pt}{10pt}} variation. Using the marked point, $k=\inblank{\blank{1.0cm}}{\gradedblank[id=q4f, type=short, points=1, answer={12}, check=numeric]{1.0cm-2pt}{10pt}}$ and
            $y=\inblank{\blank{1.6cm}}{\gradedblank[id=q4g, type=short, points=1, answer={y=12/x}, display={$12/x$}]{1.6cm-2pt}{10pt}}$.
      \item Use your equation from (b) to find $y$ when $x=8$: \inblank{\blank{1.4cm}}{\gradedblank[id=q4h, type=short, points=1, answer={1.5}, check=numeric]{1.4cm-2pt}{10pt}}. \; Would that point be
            on the drawn curve? \inblank{\blank{1.0cm}}{\gradedblank[id=q4i, type=short, points=1, answer={yes}, check=text]{1.0cm-2pt}{10pt}}
      \item Graph II has asymptotes $x=\inblank{\blank{0.8cm}}{\gradedblank[id=q4j, type=short, points=1, answer={0}, check=numeric]{0.8cm-2pt}{10pt}}$ and $y=\inblank{\blank{0.8cm}}{\gradedblank[id=q4k, type=short, points=1, answer={0}, check=numeric]{0.8cm-2pt}{10pt}}$. Graph I has none.
            Explain why a direct variation can \emph{never} have an asymptote. \vspace{1pt}\noindent\hfill\llap{\raisebox{0pt}[0pt][0pt]{\gradedblank[id=q4l, type=rubric, points=2, display={$y=kx$ is a line: every real number is in its domain, so there is no input to build a wall at, and its outputs keep growing rather than settling toward a value}]{10cm}{10pt}}}\vspace{3pt}
    \end{enumerate}

  \item \textbf{Error analysis.} The problem: \emph{$y$ varies inversely as $x$, and $y=8$ when $x=3$.
        Find $y$ when $x=12$.} A student wrote:
    \begin{center}
    \fbox{\parbox{0.80\linewidth}{\centering \vspace{2pt}
    ``Inverse, so $k=xy=(3)(8)=24$. \\
    The equation is $y=24x$. \\
    When $x=12$, \; $y=24(12)=288$.''\vspace{2pt}}}
    \end{center}
    \begin{enumerate}[label=(\alph*), leftmargin=2.1em, itemsep=5pt]
      \item Is the student's $k$ correct? \inblank{\blank{1.0cm}}{\gradedblank[id=q5a, type=short, points=1, answer={yes}, check=text]{1.0cm-2pt}{10pt}} \; So the mistake is not in finding $k$ ---
            it is in \inblank{\blank{4.0cm}}{\gradedblank[id=q5b, type=short, points=1, answer={writing the equation}, check=text]{4.0cm-2pt}{10pt}}.
      \item Test their equation on the pair they were \emph{given}: $y=24x$ at $x=3$ gives
            \inblank{\blank{1.4cm}}{\gradedblank[id=q5c, type=short, points=1, answer={72}, check=numeric]{1.4cm-2pt}{10pt}}, but it should give \inblank{\blank{1.0cm}}{\gradedblank[id=q5d, type=short, points=1, answer={8}, check=numeric]{1.0cm-2pt}{10pt}}.
      \item Write the correct equation \inblank{\blank{2.0cm}}{\gradedblank[id=q5e, type=short, points=1, answer={y=24/x}, display={$y=24/x$}]{2.0cm-2pt}{10pt}} and the correct answer \inblank{\blank{1.4cm}}{\gradedblank[id=q5f, type=short, points=1, answer={y=2}]{1.4cm-2pt}{10pt}}.
      \item There was a one-second sanity check available before any arithmetic. What is it?
            \vspace{1pt}\noindent\hfill\llap{\raisebox{0pt}[0pt][0pt]{\gradedblank[id=q5g, type=rubric, points=2, display={in an inverse variation, when $x$ goes up ($3\to12$) $y$ must go \emph{down} --- an answer of $288$ instead of $8$ is moving the wrong way}]{2.8cm}{10pt}}}\vspace{3pt}
    \end{enumerate}

  \item \textbf{The seesaw.} To balance a seesaw, the distance $d$ a child sits from the pivot varies
        inversely as their weight $w$. A $60$-pound child balances at $5$ feet from the pivot.
    \begin{enumerate}[label=(\alph*), leftmargin=2.1em, itemsep=5pt]
      \item $k=\inblank{\blank{1.4cm}}{\gradedblank[id=q6a, type=short, points=1, answer={300}, check=numeric]{1.4cm-2pt}{10pt}}$, measured in \inblank{\blank{2.6cm}}{\gradedblank[id=q6b, type=short, points=1, answer={pound-feet}, check=text]{2.6cm-2pt}{10pt}}. \; Equation: $d=\inblank{\blank{1.8cm}}{\gradedblank[id=q6c, type=short, points=1, answer={d=300/w}, display={$300/w$}]{1.8cm-2pt}{10pt}}$
      \item A $75$-pound child balances at \inblank{\blank{1.4cm}}{\gradedblank[id=q6d, type=short, points=1, answer={4}, check=numeric]{1.4cm-2pt}{10pt}} feet. A $50$-pound child balances at
            \inblank{\blank{1.4cm}}{\gradedblank[id=q6e, type=short, points=1, answer={6}, check=numeric]{1.4cm-2pt}{10pt}} feet.
      \item As $w$ gets very large, $d$ approaches \inblank{\blank{1.0cm}}{\gradedblank[id=q6f, type=short, points=1, answer={0}, check=numeric]{1.0cm-2pt}{10pt}}. Say what that horizontal asymptote
            means for a very heavy rider. \inblank{\blank{5.4cm}}{\gradedblank[id=q6g, type=rubric, points=2, display={a very heavy rider must sit almost on top of the pivot
            --- but never exactly at it}]{5.4cm-2pt}{10pt}}
      \item The vertical asymptote is at $w=\inblank{\blank{0.8cm}}{\gradedblank[id=q6h, type=short, points=1, answer={0}, check=numeric]{0.8cm-2pt}{10pt}}$. What does \emph{it} say?
            \inblank{\blank{5.4cm}}{\gradedblank[id=q6i, type=rubric, points=2, display={the lighter the rider, the further out they must sit, without bound --- a weightless
            rider could never balance the seesaw at all}]{5.4cm-2pt}{10pt}}
      \item The seesaw beam is only $8$ feet long on each side of the pivot. So the equation is only
            useful for weights of at least \inblank{\blank{1.6cm}}{\gradedblank[id=q6j, type=short, points=2, answer={37.5}, check=numeric]{1.6cm-2pt}{10pt}} pounds. Show the inequality you solved.
            \inblank{\blank{4.6cm}}{\gradedblank[id=q6k, type=short, points=1, answer={300/w<=8}, display={$300/w\le 8 \Rightarrow 300\le 8w \Rightarrow w\ge 37.5$}]{4.6cm-2pt}{10pt}}
      \item The equation cheerfully accepts $w=-60$ and returns $d=-5$. You reject it. Is that the same
            kind of rejection as an \emph{extraneous} solution from Lesson 4.6? Explain.
            \gradedlines{2}{\gradedblank[id=q6l, type=rubric, points=3, display={No. $w=-60$ genuinely satisfies $d=\frac{300}{w}$ --- the algebra has no complaint
            with it at all. It is rejected because weights and distances cannot be negative, which is the
            \emph{situation} talking. An extraneous solution never satisfied the original equation in
            the first place; it made a denominator zero.}]{\linewidth}{\gradedlinesheight}}
    \end{enumerate}
\end{enumerate}
\end{notesbox}

\begin{extensionbox}
\textbf{Part 1 --- Two variables at once.}
\begin{enumerate}[label=(\alph*), leftmargin=2.1em, itemsep=4pt]
  \item $y$ varies jointly as $x$ and $z$; $y=90$ when $x=5$ and $z=3$. \; $k=\blank{1.0cm}$, \;
        $y=\blank{1.6cm}$, \; and $y=\blank{1.2cm}$ when $x=2$ and $z=7$.
  \item $y$ varies directly as $x$ and inversely as $z$; $y=9$ when $x=6$ and $z=4$. \;
        $k=\blank{1.0cm}$, \; $y=\blank{1.6cm}$, \; and $y=\blank{1.2cm}$ when $x=10$ and $z=5$.
\end{enumerate}

\vspace{4pt}
\textbf{Part 2 --- Prove it, do not just say it.}
\begin{enumerate}[label=(\alph*), leftmargin=2.1em, itemsep=4pt]
  \item $y$ varies inversely as $x$. If $x$ is multiplied by $5$, then $y$ is multiplied by
        \blank{1.2cm}. Show it by replacing $x$ with $5x$ in $y=\frac kx$. \writeline
  \item $y$ varies jointly as $x$ and $z$. If $x$ is doubled and $z$ is halved, $y$ \blank{3.0cm}.
        Show it. \writeline
\end{enumerate}

\vspace{4pt}
\textbf{Part 3 --- Can it be both?} Ana claims a table could be a direct variation \emph{and} an
inverse variation at the same time.
\begin{enumerate}[label=(\alph*), leftmargin=2.1em, itemsep=4pt]
  \item If both were true, then $y=kx$ and $y=\frac mx$ at every row, so $kx=\frac mx$ and therefore
        $x^2=\blank{1.4cm}$.
  \item How many different values of $x$ can satisfy that? \blank{1.2cm} \; So Ana's table could
        have at most \blank{1.2cm} distinct $x$-values, and no real table of four inputs can do it.
  \item A direct variation and an inverse variation both pass through $(2,6)$. Their equations are
        $y=\blank{1.2cm}$ and $y=\blank{1.4cm}$. Where else do the two graphs cross?
        \blank{2.4cm}
\end{enumerate}
\end{extensionbox}

\begin{spiralbox}
\textbf{Coming next --- the Unit 4 Test.} That is the whole unit: what a rational function \emph{is}
and how its graph behaves (4.0), simplifying and the restrictions that come off the \emph{original}
denominator (4.1), multiplying and dividing (4.2), adding and subtracting with an LCD (4.3), the parent
$\frac1x$ and its transformations (4.4), graphing from the equation --- holes, walls, and the
horizontal asymptote by degree comparison (4.5), solving equations and throwing out the extraneous
answers (4.6), and today, the same function doing an honest day's work. The practice test in your
packet mirrors the real one in format and ideas with different numbers. One habit is worth more than
any formula on it: \textbf{factor first, restrict early, and check what the answer is supposed to
mean.}
\end{spiralbox}

\end{document}
